\section{Degree-floor geometry and sharp-transfer tools}
\label{sec:sparse-extension}

This section collects the regularized geometry and sharp-transfer tools
used in \cref{thm:regularized}. The earlier unregularized encoder
rejects a graph with any zero degree; that convention cannot be
retained for the sparse result. The degree-plus-spectrum construction
already present in the project avoids this algebraic difficulty.
The tools below require no logarithmic degree lower bound.
Their row-tail, local full-to-LOO, and fitted-parameter interfaces,
left open in revision 2, are now proved in
\cref{sec:sparse-restored} for the explicitly regularized algorithm.

\subsection{The existing degree-plus-spectrum action remains valid}

Write $\mathbf d=A\one$ for the raw degree vector. For any strictly
positive diagonal matrix $W$, retain
\[
 WAWU=U\Lambda,\qquad U^\top U=I,\qquad
 X=[\one,WU],\qquad Y=[\mathbf d,W^{-1}U\Lambda].
\]
Then the exact identities
\begin{equation}\label{eq:sparse-old-action}
 AX=Y,\qquad YX^\dagger=A\Pi_X,\qquad A\Pi_X\one=\mathbf d
\end{equation}
hold for every graph, including graphs with zero degrees. Here $\Pi_X$
is the Euclidean orthogonal projector onto the column space of $X$;
no full-column-rank assumption is necessary. The earlier choice
$W_{ii}=\min(1,2\bar D/D_i)$ for $D_i>0$, with $W_{ii}=1$ for
$D_i=0$ and $\bar D=n^{-1}\sum_iD_i$, supplies a positive reweighting
when $\bar D>0$. The all-zero graph can be assigned a fixed output.
The block-count error is exactly $-A(I-\Pi_X)Z$.

These facts need no logarithmic expected-degree assumption. They are
computational identities, however: their validity does not give the
sharp tail of the remaining random block-count error. In DC-SBM a
normalization must also preserve the effective likelihood coefficients,
as opposed to preserving only a population mean.

\subsection{Degree flooring preserves the relevant DC geometry}

Use the completed mean from \eqref{eq:mean-model} and put
$D_0=\diag(\Omega\one)$. Its diagonal is between fixed positive
multiples of $d$. Choose a fixed $a_0>0$ below its minimum structural
multiplier and define
\[
 G_i=\max(D_i,a_0d),\qquad G=\diag(G_i),\qquad
 S_G=G^{-1/2}AG^{-1/2}.
\]
At the deterministic population degrees the floor is inactive.
In particular $D_0^{-1}\Theta Z$ is block constant; adding a multiple
of $d$ to every degree would not have this property for heterogeneous
theta.

\begin{proposition}[Sparse degree-floor geometry]\label{prop:sparse-floor}
Under the fixed structural conditions \eqref{eq:mean-model}--%
\eqref{eq:structural}, suppose $d=np\to\infty$ and $np^2\le C_1$.
For every fixed $r>0$, with probability at least $1-C_rn^{-r}$,
\begin{align}
 \norm{G^{-1/2}(A-\Omega)G^{-1/2}}&\le C_r d^{-1/2},
 &\sum_iG_i&\le C_rnd,\label{eq:floor-noise}\\
 \sum_i\frac{(D_i-(D_0)_{ii})^2}{G_i}
 +\sum_i\frac{(G_i-(D_0)_{ii})^2}{G_i}
 &\le C_rn.\label{eq:floor-degree-energy}
\end{align}
Let $U$ contain the $K$ eigenvectors of $S_G$ with largest absolute
eigenvalues, let $\Pi=UU^\top$, and let
$H=G^{-1/2}\Pi G^{1/2}$. For all sufficiently large $n$, on the same
event the retained eigenvalues have a fixed separation from zero and
\begin{equation}\label{eq:floor-block-residual}
 \norm{(I-\Pi)G^{1/2}Z}_F\le C_r\sqrt n.
\end{equation}
For every block-constant $f$ with $\norm f_\infty\le C_f$,
\begin{equation}\label{eq:floor-weighted-errors}
 \norm{(I-H)f}_G\le C_r\sqrt n,\qquad
 \norm{(I-H)f}_2^2\le C_rn/d,
 \quad \norm v_G^2=v^\top Gv.
\end{equation}
No lower bound proportional to $\log n$ is imposed.
\end{proposition}

\begin{proof}
Let $d_*=\max_{i\ne j}n\,\E A_{ij}\asymp d$, and set
$W_{h,ii}=\sqrt{\min(1,2d_*/D_i)}$, with value one at zero degrees.
The changed set $\mathcal S=\{i:D_i>2d_*\}$ has at most $10n/d_*$
vertices except on an event of probability $e^{-cn}$ for large $n$.
Here is a direct verification of the probability strength. If
$m=\lfloor10n/d_*\rfloor+1\le n$ vertices all exceed that degree
threshold, some set $T$ of size $m$ has
$\sum_{i\in T}D_i>2md_*$. This sum is an independent-edge sum with
coefficients in $\{1,2\}$ and mean at most $md_*$. At
$t=(\log2)/2$ its Bernoulli exponential transform gives a tail at most
\[
 \exp\{-(\log2-1/2)md_*\}.
\]
Indeed $(e^{tc}-1)/c\le(e^{2t}-1)/2=1/2$ for $c=1,2$.
The logarithm of the number of sets is
$\log\binom nm\le m\log(en/m)=o(md_*)$, because $d_*\to\infty$.
This proves the asserted $e^{-cn}$ bound. If $m>n$ there is no
exceptional set to consider.

The squared Euclidean norm of every row of $W_hAW_h$ is at most
$W_{h,ii}^2D_i\le2d_*$. Apply the row-norm form of the regularized
adjacency theorem of Le, Levina and Vershynin
\cite[Theorem 2.1 and Remark 2.2]{llv2017}. For a fixed sufficiently
large choice of its probability parameter, it gives
\[
 \norm{W_hAW_h-\E A}\le C_r\sqrt d
\]
with failure probability at most $C_rn^{-r}$. The weights are allowed
to depend on the graph. Also
\[
 \norm{W_h\Omega W_h-\Omega}
 \le Cp\sqrt{n|\mathcal S|}\le C\sqrt d;
\]
the matrix on the left is supported on rows or columns in
$\mathcal S$, so this follows already from its Frobenius norm.
The completed-mean diagonal has norm $O(p)$. These observations imply
$\norm{W_h(A-\Omega)W_h}\le C_r\sqrt d$.

There is a diagonal $F$ with bounded norm, depending only on the
structural constants and $a_0$, such that
\[
 G^{-1/2}=(2d_*)^{-1/2}FW_h.
\]
For $D_i>2d_*$ its diagonal entry is one, and otherwise it is
$\sqrt{2d_*/G_i}$. Conjugating the preceding bound proves the first
part of \eqref{eq:floor-noise}. Independently, the Bernoulli transform
of the total number of edges shows
$\sum_iD_i\le Cnd$ except with probability $e^{-cnd}$.
Since $G_i\le D_i+a_0d$, this proves the second part.

Apply the normalized noise to $G^{1/2}\one$:
\[
 \norm{G^{-1/2}(A-\Omega)\one}_2
 \le C_rd^{-1/2}\norm{G^{1/2}\one}_2\le C_r\sqrt n.
\]
This proves the first sum in \eqref{eq:floor-degree-energy}.
Because $a_0d<(D_0)_{ii}$, flooring can only reduce the distance to
$(D_0)_{ii}$:
$|G_i-(D_0)_{ii}|\le|D_i-(D_0)_{ii}|$.
This proves the second sum without any separate degree-wise union bound.

The nonzero eigenvalues of $G^{-1/2}\Omega G^{-1/2}$ have a fixed
gap from zero. To verify this without a maximum-degree assumption, its
factor $V=G^{-1/2}\Theta Z$ has diagonal Gram entries
\[
 \frac{cn}{d}
 \le
 \frac{(\sum_{i:z_i=a}\theta_i)^2}{\sum_{i:z_i=a}G_i}
 \le \sum_{i:z_i=a}\frac{\theta_i^2}{G_i}
 \le\frac{Cn}{d}.
\]
The lower bound uses comparable group sizes and the total-degree
upper bound; the upper bound uses the floor. Thus the nonzero spectrum
is that of $p(V^\top V)^{1/2}C(V^\top V)^{1/2}$ and is separated
from zero by a fixed amount. Spectral perturbation and
\eqref{eq:floor-noise} give projector error $O(d^{-1/2})$.

Now $G^{-1/2}D_0Z$ belongs to the population signal space, whereas
\[
 \norm{G^{1/2}Z-G^{-1/2}D_0Z}_F^2
 =\sum_i\frac{(G_i-(D_0)_{ii})^2}{G_i}\le C_rn.
\]
Moreover $\norm{G^{1/2}Z}_F^2=\sum_iG_i\le C_rnd$.
Combining these facts with the projector error proves
\eqref{eq:floor-block-residual}. Since $H$ is the $G$-orthogonal
projector corresponding to $\Pi$, the first bound in
\eqref{eq:floor-weighted-errors} follows. The second uses $G_i\ge a_0d$.
\end{proof}

The scale $d$ in this proof need not be an algorithm input. One may use
a sufficiently small fixed multiple of $\bar D$ for the floor. The
same total-edge Bernoulli transform also gives the lower bound
$\sum_iD_i\ge cnd$ with failure probability $e^{-c'nd}$.
On this two-sided total-edge event, the threshold lies between fixed positive
multiples of $d$ and below every $(D_0)_{ii}$. Every deterministic step
above then applies. The formula is defined even at isolated vertices.
Deleting a vertex and recomputing degrees gives the same statements on
the remaining graph, with its completed mean and block masses. Choosing
the probability parameter larger also permits a union over all such
deleted graphs. This does not yet compare their scores to the full
graph score.

\subsection{The DC tilted mean and quadratic error already close}

Work on a deleted graph so that all coordinates below index actual
neighbors of a separate tested vertex. For the bounded block-constant
linear candidate contrast used in the DC profile reduction,
the exact tilted Bernoulli probabilities have the decomposition
\begin{equation}\label{eq:floor-tilt-decomposition}
 r=D_0h+\epsilon,\qquad
 h\text{ block constant},\quad
 \norm h_\infty\le C/n,\quad \norm\epsilon_\infty\le Cp^2.
\end{equation}
For example the leading term of the exact tilt is
$\theta_i\theta_jpC_{a,z_j}e^{-s w_{z_j}}$, while
$(D_0)_{jj}=\theta_jp\sum_k\theta_k C_{z_j,z_k}$.
Dividing the former by the latter gives the required block-constant
$h_j$. The denominator in the exact Bernoulli tilt contributes
$O(p^2)$ uniformly. This uses the original Bernoulli law throughout.
The target coordinate is absent; artificially appending $r_i=0$ would
instead create an $O(p)$ error at that one coordinate.

For bounded block-constant $f$, put $e=(H-I)f$. The
$G$-orthogonality of $H$ gives the exact cancellation
\begin{align}
 r^\top e
 ={}&-\langle(I-H)h,(I-H)f\rangle_G
       +((D_0-G)h)^\top e+\epsilon^\top e.
       \label{eq:floor-tilted-cancellation}
\end{align}
The first term is $O(1)$ by \eqref{eq:floor-block-residual}, since
the block coefficients of $h$ are $O(1/n)$. For the second,
Cauchy--Schwarz and \eqref{eq:floor-degree-energy} give
\[
 |((D_0-G)h)^\top e|
 \le \norm h_\infty
 \left(\sum_j\frac{((D_0)_{jj}-G_j)^2}{G_j}\right)^{1/2}
 \norm e_G
 \le C.
\]
The last term is at most
$Cp^2\sqrt n\,\sqrt{n/d}\le Cnp^2$. Also $r_j\le Cp$, so
\begin{equation}\label{eq:floor-tilt-moments}
 |r^\top e|\le C(1+np^2),\qquad
 \sum_jr_je_j^2\le Cp\norm e_2^2\le C.
\end{equation}
Thus the first-order bias cancellation and bounded quadratic error do
not require logarithmic degree. The replacement for a uniform
coefficient bound is supplied by
\cref{restore:tail-energy-lemma,restore:tail-hit-corollary}.

\subsection{Exact preservation of the raw degree vector}

Let $S_GU=U\Lambda$, put $H=G^{-1/2}UU^\top G^{1/2}$, and define
\begin{equation}\label{eq:floor-degree-corrected-action}
 R_0=G^{1/2}U\Lambda U^\top G^{1/2}=AH,\qquad
 R=R_0+(\mathbf d-R_0\one)\frac{\one^\top}{n}.
\end{equation}
All quantities in this formula are computable from the retained degrees
and eigenpairs. Exactly $R\one=\mathbf d$. For every bounded
block-constant likelihood vector $w$, with $\bar w=n^{-1}\one^\top w$,
\begin{equation}\label{eq:floor-effective-coefficients}
 Rw=A\{H(w-\bar w\one)+\bar w\one\}.
\end{equation}
Hence the coefficient error is precisely $(H-I)(w-\bar w\one)$,
to which \eqref{eq:floor-tilt-moments} applies on a deleted graph.
This construction avoids orthogonalizing a possibly tiny residual
constant vector. We do not assert that $R$ is an orthogonal-projection
reconstruction.

A hollow version subtracts $R_{ii}$ from entry $(i,i)$ and adds
$R_{ii}/(n-1)$ to each other entry of row $i$. This preserves raw
degrees and changes a bounded row contrast by at most
$2\norm w_\infty |R_{ii}|$. A sufficiently strong local bound for
$R_{ii}$ must therefore be proved before this hollowing can be inserted
into a sparse algorithmic theorem; it does not follow from aggregate
geometry alone.

\subsection{Two sharp transfers that do not assume every coefficient is bounded}

The next two lemmas give concrete alternatives to a global
row-incoherence event. They concern a row-independent, training-measurable
score, just as in \cref{lem:sharp}.

\begin{lemma}[Exponential-moment coefficient control]\label{lem:sparse-exp-transfer}
Assume the oracle hypotheses of \cref{lem:sharp}. Let
$\widehat T=T+b+\sum_je_jX_j$, with $(b,e)$ measurable from independent
training data. On a measurable training event suppose
\[
 \left|b+\sum_jr_je_j\right|\le B,\qquad
 \sum_jr_je_j^2e^{|e_j|}\le V,
\]
where $B,V$ are fixed. Then, for all real $x$ and all sufficiently
large $d$,
\[
 \PP(\widehat T\le x\mid\text{training})
 \le C e^{s_*x+B+V/2}d^{-1/2}e^{-I}.
\]
The constant is independent of the largest coefficient.
\end{lemma}

\begin{proof}
Taylor's inequality gives, under the original oracle tilt $Q$,
\[
 \log\E e^{-s_*\widehat T}
 \le-I-s_*\left(b+\sum_jr_je_j\right)
       +\frac12\sum_jr_je_j^2e^{|e_j|}
 \le-I+B+V/2.
\]
Tilt again by $e^{-s_*\widehat T}$. The new probabilities satisfy
\[
 \widehat r_j
 =\frac{r_je^{-s_*e_j}}{1-r_j+r_je^{-s_*e_j}}
 \le r_je^{s_*|e_j|}.
\]
For $|e_j|\le1$, this bounds the new centered third absolute moments
by $C_W\sum_jr_j=O(d)$. For $|e_j|>1$, use $s_*\le1-\tau$ and
\[
 e^{s_*|u|}(W+|u|)^3
 \le C_{W,\tau}u^2e^{|u|}\qquad (|u|>1).
\]
The remaining third moments are therefore $O(V)$. The upper variance
bound follows from the same estimates. There is no restriction that
every $e_j$ be bounded.

For the lower variance, restrict to $|e_j|\le1$. The original variance
discarded by this restriction is at most $W^2V$, and the weighted
$\ell_2$ norm of the remaining perturbation is at most $\sqrt V$.
Thus the square root of the original-weight variance of the remaining
$(w_j+e_j)$ is at least $\sqrt{vd-W^2V}-\sqrt V$.
On this restricted set, the new and old Bernoulli variances differ by
only fixed factors. The new total variance is therefore at least $cd$
for large $d$. Berry--Esseen gives a $C/\sqrt d$ bound for every unit
interval under the new tilt. Summing those intervals, using the first
displayed exponential-moment bound, proves the assertion as in
\cref{lem:sharp}.
\end{proof}

\begin{lemma}[Paying exceptional coefficients only when the row hits them]
\label{lem:sparse-truncated-transfer}
Assume the oracle hypotheses of \cref{lem:sharp}, and write $p_j$
for the un-tilted row probabilities. Suppose $p_j\le C_pr_j$, as
holds here for bounded oracle likelihood coefficients. On an independent
training event $\mathcal E$, suppose
\[
 \left|b+\sum_jr_je_j\right|\le B,\qquad \sum_jr_je_j^2\le V.
\]
Here $B,V$ are fixed constants. For any fixed $M>0$ define
\[
 h_M=\sum_{j:|e_j|>M}p_j,\qquad
 \eta_M=\sum_jr_je_j^2\ind{|e_j|>M}.
\]
Then, for every fixed bounded $x$,
\begin{equation}\label{eq:sparse-truncated-transfer}
 \PP(\widehat T\le x)
 \le C_Md^{-1/2}e^{-I}
       +\E[\ind{\mathcal E}h_M]
       +\PP(\mathcal E^c).
\end{equation}
Since $h_M\le C_p\eta_M/M^2$, an expected truncated quadratic error of order
$d^{-1/2}e^{-I}$ suffices. A bound $Ce^{-Hd}$ is a convenient
sufficient condition whenever $I\le C_Id$ and $H>C_I$.
The expected hit-mass condition on $h_M$ is weaker and is itself enough.
\end{lemma}

\begin{proof}
For analysis only, replace $e_j$ by
$e'_j=e_j\ind{|e_j|\le M}$. The truncated coefficients satisfy
\[
 \left|b+\sum_jr_je'_j\right|
 \le B+\eta_M/M\le B+V/M,\qquad
 \sum_jr_j(e'_j)^2\le V.
\]
Hence \cref{lem:sharp} bounds their score error by
$C_Md^{-1/2}e^{-I}$ on $\mathcal E$.
The original and truncated random scores are identical whenever
$X_j=0$ for every $j$ with $|e_j|>M$. Conditional on training, the
probability that this fails is at most
\[
 \sum_{j:|e_j|>M}p_j
 =h_M\le \frac{C_p}{M^2}\eta_M.
\]
Integrate this bound on $\mathcal E$ and pay at most
$\PP(\mathcal E^c)$ outside it. No coefficient truncation is an
additional input or computation required of the implemented decoder.
\end{proof}

For the degree-floor projection there is a concrete reduction to a
weighted row count. Define
\[
 \zeta_j=\sqrt{\frac{nd}{G_j}}\norm{U_{j\cdot}}_2.
\]
If $e=(H-I)f$, $\norm f_\infty\le C_f$ and
$\sum_jG_j\le C_Gnd$, then deterministically
\[
 |e_j|
 \le \frac{\norm{U_{j\cdot}}_2}{\sqrt{G_j}}
        \norm{U^\top G^{1/2}f}_2+C_f
 \le C_f\sqrt{C_G}\,\zeta_j+C_f.
\]
Choosing $M=C_f\sqrt{C_G}L+C_f$ thus gives
$\{|e_j|>M\}\subseteq\{\zeta_j>L\}$. Since $p_j\le Cd/n$,
the sufficient expected-count estimate
\begin{equation}\label{eq:sparse-row-count-target}
 \E\!\left[\ind{\mathcal E}
       \sum_j\ind{\zeta_j>L}\right]\le Cne^{-Hd},
 \qquad H>C_I,
\end{equation}
would imply
$\E[\ind{\mathcal E}h_M]\le Cd e^{-Hd}
=o(d^{-1/2}e^{-I})$ whenever $I\le C_Id$.
Equation \eqref{eq:sparse-row-count-target} is not a consequence of
aggregate geometry alone. It is proved by the same-threshold
truncated-energy argument in \cref{restore:tail-energy-lemma}
and the floor lower bound.

\subsection{The interfaces completed in the sparse restoration}

The preceding propositions prove more than spectral consistency:
the weighted projection has bounded tilted mean and variance, and the
sharp tail permits exceptional coefficients whose expected hit mass is
small. Three additional estimates turn those tools into the actual
algorithm theorem. First, \cref{restore:tail-energy-lemma} proves the
needed truncated energy on the full and recomputed deleted graphs.
Second, \cref{restore:loo-scalar,restore:loo-corrected} compare the
actual full-graph score to the independent deleted score, including
the empirical floor, exact degree correction, and hollowing, at
$O(1)$ cost using row bounds only on incident coordinates.
Third, \cref{restore:fit} controls initialization and fitted likelihood
contrasts by aggregate weighted errors, while retaining exceptional
vertices. The invariant-basin argument covers finite hard refits
without a union over their number.

For $d=o(\log n)$ the already proved $n^{-r}$ geometry failure is
negligible relative to $d^{-1/2}e^{-C_Id}$ for any fixed $r>0$.
This explains why the regularized norm event needs no stronger
probability in that range. Nor should one declare failure of the entire
graph whenever any vertex has an unusually small degree. The
truncated-score comparison shows how local exceptional coordinates can
instead be paid when the tested row encounters them.

Zhang's spectral clustering proof
\cite[Theorem 3 and Corollary 1, p.~7326]{zhang} establishes exponentially
small expected truncated row energy and exceptional-row counts for
the trimmed adjacency eigenspace. That result provides a concrete
model for \eqref{eq:sparse-row-count-target}; its eigenspace and
normalization differ from the floor-normalized DC construction here.
Its final leading-exponent conclusion does not, by itself, imply
the bounded perturbation cost and sharp prefactor required here.
The intended nontrivial sparse target on the present fixed-separation,
bounded-theta class is $d=np\to\infty$, together with the stated
Bernoulli remainder condition. A vanishing signal gap or degree
parameters approaching zero requires additional quantitative work.
The completed proof in \cref{sec:sparse-restored} establishes each
of these estimates. It does not replace them by an assumption that
the desired constant risk ratio already holds.

